\documentclass{leveldoc}
\DocUses{foundations,first_bend,geiger_marsden}
\DocFingerprint{b9ed631a76c036fc}

\begin{document}

\section{What the level asks}\label{sec:asks}

The game's text:
\leveldescription

Two particles, each with $q = \val{q}$, $m = \val{m}$ and kinetic energy $T_0 = \val{Tzero}$,
leave the nodes $(0, 13)$ and $(0, 7)$ along $+x$. Each has its own detector, and neither
detector sets a condition on direction or on energy: the first entry solves that shot. The
upper detector is the rectangle $27 \le x \le 30$, $15 \le y \le 19$; the lower one is
$27 \le x \le 30$, $1 \le y \le 5$. The player places up to three charges, of either sign
and of the magnitudes in the level file, on any node of the grid. There are no fixed
elements. A flight ends when the particle touches a charge, leaves the world or runs out of
time (\xsecref{foundations}{sec:board}); $t_{\max} = \val{tmax}$. The level is solved when
both shots arrive.

The reference solution is a single charge $Q = \val{Q}$ at $\val{rC}$. It is repulsive, and
it sits between the two rays, four cells below the upper launch and two cells above the
lower one. It pushes the rays apart. The upper particle, on the wider offset, turns upward
into the high detector; the lower particle, on the narrower offset, turns downward more
sharply into the low detector. Each flight is the orbit about one charge, the closed form
of \levelref{first_bend}, evaluated at the detector's left face. The time is Simpson's rule
on that orbit.

\section{Setup}\label{sec:setup}

The level file, field by field (the appendix reproduces it):
\setuptable

In symbols: $c = \val{c}$, each particle has radius $a = 0$, and the charge has radius
$R = \val{R}$. The coupling of the reference charge is $\kappa = qQ = \val{kappa}$. Both
launches have
\[
  \gamma_0 = 1 + \frac{T_0}{mc^2} = \val{gammaZero}, \qquad
  \abs{\vect{p}_0} = \sqrt{\frac{T_0^2}{c^2} + 2mT_0} = \val{pZero},
\]
and the direction in the file is $(1, 0)$. The world is $-1 \le x \le 31$, $-1 \le y \le 21$.
The two launches and the two detectors are mirror images through the line $y = 10$
(§\ref{sec:mirror}). The reference charge is one cell below that line.

\begin{figure}[htbp]
\centering
\begin{tikzpicture}
\begin{axis}[axis equal image, width=0.92\linewidth, xmin=-1.6, xmax=31.6, ymin=-1.6, ymax=21.6,
  xtick={0,5,...,30}, ytick={0,5,...,20}, tick label style={font=\small},
  xlabel={$x$}, ylabel={$y$}, label style={font=\small}]
  \draw[black!8] (0,0) grid (30,20);
  \draw[black!45, dashed] (-1,-1) rectangle (31,21);
  \draw[black!35, dashed] (0,10) -- (30,10);
  \fill[green!50!black, fill opacity=0.22] (27,15) rectangle (30,19);
  \fill[green!50!black, fill opacity=0.22] (27,1) rectangle (30,5);
  \draw[green!40!black] (27,15) rectangle (30,19);
  \draw[green!40!black] (27,1) rectangle (30,5);
  \addplot[very thick, blue!70!black] table {pathup.dat};
  \addplot[very thick, orange!80!black] table {pathdown.dat};
  \fill[red!70!black] (24,9) circle[radius=0.28];
  \node[font=\small, red!60!black, anchor=south] at (24,9.35) {$Q$};
  \draw[-{Stealth[length=6pt]}, thick] (0,13) -- (2.1,13);
  \draw[-{Stealth[length=6pt]}, thick] (0,7) -- (2.1,7);
  \fill (0,13) circle[radius=0.12];
  \fill (0,7) circle[radius=0.12];
  \node[font=\small, anchor=south west] at (2.2,13.05) {upper};
  \node[font=\small, anchor=north west] at (2.2,6.9) {lower};
  \node[font=\small, green!30!black, anchor=west] at (27.15,17) {upper};
  \node[font=\small, green!30!black, anchor=west] at (27.15,3) {lower};
\end{axis}
\end{tikzpicture}
\caption{The reference flights. The upper particle (blue) and the lower particle (orange)
each enter its detector through the face $x = 27$. The dashed horizontal line is the mirror
$y = 10$ of §\ref{sec:mirror}; the charge sits one cell below it.}
\label{fig:board}
\end{figure}

\section{Principles and model}\label{sec:model}

\begin{principles}
  \item The field of a charged sphere and the force on a non-polarizable particle,
    \xeqref{foundations}{eq:coulomb} and \xeqref{foundations}{eq:sphereforce}.
  \item The relativistic equation of motion \xeqref{foundations}{eq:eom}, the launch
    \xeqref{foundations}{eq:launch}, the energy \xeqref{foundations}{eq:energy} and the
    angular momentum \xeqref{foundations}{eq:angmom}. With one charge, the orbit
    \xeqref{first_bend}{eq:general}.
  \item The board's rules (\xsecref{foundations}{sec:board}): a level with several shots is
    solved when every shot arrives.
  \item Newton's method and Simpson's rule (\xsecref{foundations}{sec:numerics}).
  \item The Newtonian scattering angle of one charge for a particle that comes from
    infinity, \xeqref{geiger_marsden}{eq:chi-tan}. It estimates the turn; the arrival
    compared below is the orbit at the face.
\end{principles}

With one charge at $\vect{r}_\mathrm{c}$ the state of either particle obeys
\begin{equation}\label{eq:motion}
  \deriv{\vect{x}}{t} = \frac{\vect{p}}{\gamma m}, \qquad
  \deriv{\vect{p}}{t} = \kappa\,\frac{\vect{x}-\vect{r}_\mathrm{c}}{\abs{\vect{x}-\vect{r}_\mathrm{c}}^3},
  \qquad \gamma = \sqrt{1 + \frac{\abs{\vect{p}}^2}{m^2c^2}},
\end{equation}
from its own launch, until the first event. The motion stays in the plane
(\xsecref{foundations}{sec:plane}). The two flights do not interact: each particle flies
alone.

\begin{assumption}{The charge is held in place}\label{as:fixed}
The charge does not recoil. That is a rule of the game (\xsecref{foundations}{sec:radiation}).
\end{assumption}

\begin{assumption}{No radiation}\label{as:radiation}
The radiated energy \xeqref{foundations}{eq:larmor} is not taken from the particle. Along
the two reference flights it is $\val{EradUp}$ and $\val{EradDown}$, fractions
$\val{EradFracUp}$ and $\val{EradFracDown}$ of $T_0$ (§\ref{sec:hand}). The game makes the
same approximation and reports the same integral (\physref{8}).
\end{assumption}

\begin{assumption}{A point particle}\label{as:point}
$a = 0$, and the particle does not polarize, so \xeqref{foundations}{eq:sphereforce} is
exact. The closest approaches are $\val{rminUp}$ and $\val{rminDown}$, both far outside
the sphere of radius $R$.
\end{assumption}

\begin{assumption}{Where the scattering formula applies}\label{as:formula}
The formula \xeqref{geiger_marsden}{eq:chi-tan} is Newtonian, has a single charge, and
takes the kinetic energy and the impact parameter at infinity. Here $c$ is finite and each
particle is launched a finite distance from the charge. At launch the potential energy is
the fraction $\val{potUp}$ of $T_0$ on the upper shot and $\val{potDown}$ on the lower,
while $\Gamma$ falls short of $1$ by $\val{gGapUp}$ and $\val{gGapDown}$. Continued to
infinity, the relativistic orbit turns by $\val{chiGapUp}^\circ$ and
$\val{chiGapDown}^\circ$ less than the formula (§\ref{sec:formula}). The comparison uses
the orbit at the face, which assumes none of this.
\end{assumption}

\section{Derivation}\label{sec:derivation}

\subsection{The orbit of one shot}\label{sec:orbit}

Take polar coordinates about the charge,
$\vect{x} - \vect{r}_\mathrm{c} = r(\cos\varphi, \sin\varphi)$, and let $u = 1/r$. Level~1
derives the orbit. With $\sigma = \operatorname{sign} L_z$, $L = \abs{L_z}$,
$\theta = \sigma(\varphi - \varphi_0)$ and $u'_0 = -p_{r,0}/L$,
\begin{equation}\label{eq:orbit}
  u(\theta) = B + (u_0 - B)\cos\Gamma\theta + \frac{u'_0}{\Gamma}\sin\Gamma\theta,
\end{equation}
where $\Gamma^2 = 1 - \kappa^2/(c^2 L^2)$ and $B = -\kappa W/(c^2 L^2 \Gamma^2)$, and $W$ is
the energy \xeqref{first_bend}{eq:energy}. This is \xeqref{first_bend}{eq:general}.
Differentiating it gives the slope the face condition needs,
\begin{equation}\label{eq:uprime}
  u'(\theta) = -\Gamma(u_0 - B)\sin\Gamma\theta + u'_0\cos\Gamma\theta.
\end{equation}
The time and the momentum along the same orbit are \xeqref{first_bend}{eq:dtdth} and
\xeqref{first_bend}{eq:momentum}.

Both launches point along $+x$, so $L = b\,\abs{\vect{p}_0}$ with
$b = \abs{y_\mathrm{c} - y_0}$ the vertical offset, and
$p_{r,0} = \abs{\vect{p}_0}\,(x_0 - x_\mathrm{c})/r_0$. The launch is not an apsis:
$u'_0 \ne 0$. For the reference charge the upper launch has $b = \val{bUp}$ and
$L_z < 0$, and the lower launch has $b = \val{bDown}$ and $L_z > 0$. Repulsion pushes the
upper particle up and the lower particle down.

\subsection{The face}\label{sec:face}

The position on the orbit is
\begin{equation}\label{eq:position}
  x(\theta) = x_\mathrm{c} + \frac{\cos\varphi}{u(\theta)}, \qquad
  y(\theta) = y_\mathrm{c} + \frac{\sin\varphi}{u(\theta)}, \qquad
  \varphi = \varphi_0 + \sigma\theta.
\end{equation}
The detector begins at the face $x = x_D = \val{xFace}$. While $u > 0$,
\begin{equation}\label{eq:dx}
  \deriv{x}{\theta} = -\frac{\sigma\sin\varphi}{u} - \frac{\cos\varphi}{u^2}\,u',
\end{equation}
and the same differentiation gives $\dd y/\dd\theta$. The crossing of the face is the root
of
\begin{equation}\label{eq:face}
  f(\theta) = x(\theta) - x_D = 0.
\end{equation}

Which boundary comes first is read off the orbit, not assumed. Walk $\theta$ from $0$ in
steps of $10^{-3}$, up to $\theta = 3$. On the chord from $\theta_k$ to $\theta_{k+1}$, in
the board's order (\xsecref{foundations}{sec:board}): if $r(\theta_{k+1}) \le R$, the
flight collides; if the point at $\theta_{k+1}$ lies outside the world, the flight leaves,
through the side the chord reaches first; if that point lies inside the detector, the
flight arrives, through the side the chord reaches first. The side is the one whose value
the chord attains at the smallest fraction of its length. Both reference shots meet the
face $x = x_D$ first. The root of \eqref{eq:face} is then polished by Newton's method
\xeqref{foundations}{eq:newton}, started at the middle of that chord,
$\theta_k + 5\times 10^{-4}$, with the derivative \eqref{eq:dx}.

The closest approach on the orbit is the unique root of $u' = 0$ \eqref{eq:uprime}. On
both reference flights that root lies before the face, and $r$ there is more than ten
times $R$ (§\ref{sec:hand}), so the sphere is never reached.

\subsection{The time and the radiated energy}\label{sec:time}

Level~1 evaluates
\[
  t(\theta) = \int_0^\theta \deriv{t}{\theta'}\,\dd\theta', \qquad
  \deriv{t}{\theta} = \frac{W - \kappa u}{c^2 L u^2},
\]
in closed form, because its launch is an apsis and $u = A\cos\Gamma\theta + B$
(\xeqref{first_bend}{eq:time}, \xeqref{first_bend}{eq:i1}, \xeqref{first_bend}{eq:i2}).
The same antiderivatives apply here once \eqref{eq:orbit} is written as a phase-shifted
cosine: $u'_0 \ne 0$ only moves the origin of that cosine. The number compared below is
Simpson's rule \xeqref{foundations}{eq:simpson} on $\dd t/\dd\theta$, with $\val{nSimp}$
panels. Halving the panel count estimates the error as $(S_n - S_{2n})/15$.

The radiated energy is the same rule applied to the power. The force is
$\vect{F} = \kappa(\vect{x} - \vect{r}_\mathrm{c})/r^3$, the acceleration is
\xeqref{first_bend}{eq:acceleration}, and the power is \xeqref{foundations}{eq:larmor}.
The integral uses $8000$ panels; the uncertainty is the relative step-doubling error.

\subsection{What the scattering formula estimates}\label{sec:formula}

For a repulsive charge and a particle that comes from infinity with kinetic energy $T$ and
impact parameter $b$, Level~3 derives
\begin{equation}\label{eq:chi}
  \tan\frac{\chi}{2} = \frac{\kappa}{2Tb}
\end{equation}
(\xeqref{geiger_marsden}{eq:chi-tan}). Read with $T = T_0$ and with the geometric offset
$b$, it gives $\val{chiFormUp}^\circ$ for the upper shot and $\val{chiFormDown}^\circ$ for
the lower. The relativistic orbit, continued until $u$ crosses zero on both sides of the
launch, turns by $\val{chiRelUp}^\circ$ and $\val{chiRelDown}^\circ$ between its
asymptotes. \ref{as:formula} records the difference.

The momentum at the face makes an angle $\val{angUp}^\circ$ with the launch direction on
the upper shot and $\val{angDown}^\circ$ on the lower. Those angles and the
asymptote-to-asymptote turns differ by $\val{angLeftUp}^\circ$ and $\val{angLeftDown}^\circ$.
The face lies three cells past the charge, so each orbit is cut off before the outgoing
asymptote, and each launch sits off the incoming asymptote. The formula explains which ray
turns more. The arrival is \eqref{eq:position} at the root of \eqref{eq:face}.

\subsection{The mirror}\label{sec:mirror}

The reflection $\mathcal{M}:(x, y) \mapsto (x,\, 20 - y)$ fixes the line $y = 10$ and
exchanges the two launches, the two detectors and the two edges $y = -1$ and $y = 21$. The
Coulomb field of a charge on that line has $E_x$ even and $E_y$ odd under $\mathcal{M}$,
and both launches have $p_y = 0$. The equation \eqref{eq:motion} therefore sends mirror
initial states to mirror flights: the flight times are equal, the momenta have opposite
$p_y$, and one particle enters its detector if and only if the other enters the image
detector. A charge off the line has no such duty. The reference charge is at $y = 9$, and
the two offsets are $\val{bUp}$ and $\val{bDown}$.

\section{The reference solution by hand}\label{sec:hand}

\begin{procedure}
  \item\label{st:launch} \emph{The launch}, the same for both shots
    (\xeqref{foundations}{eq:launch}): $\gamma_0 = \val{gammaZero}$ and
    $\abs{\vect{p}_0} = \val{pZero}$.
  \item\label{st:constants} \emph{The upper orbit} \eqref{eq:orbit}. The offset is
    $b = \val{bUp}$, so $L = b\,\abs{\vect{p}_0} = \val{LUp}$ and
    $r_0 = \val{rUp}$. Then $W = \val{WUp}$, $\Gamma = \val{GammaUp}$,
    $B = \val{BUp}$, $u_0 = \val{uUp}$ and $u'_0 = \val{duUp}$.
  \item\label{st:scan} \emph{The first event}, by the walk of §\ref{sec:face}. The upper
    shot meets $x = x_D$ before a collision, a world edge or any other face of its
    detector. Five points of the orbit, from the same formula:
    \begin{center}\small\val{samples}\end{center}
    The root of $u' = 0$ lies at $\theta = \val{thMinUp}$, where $r = \val{rminUp}$.
    Sampling $r$ at $401$ angles reproduces it to $\val{rminGapUp}$.
  \item\label{st:newton} \emph{The crossing.} Newton's method on \eqref{eq:face}, with
    \eqref{eq:dx}, from the middle of the bracketing chord:
    \begin{center}\small\val{newton}\end{center}
    After $\val{newtonSteps}$ corrections, $\theta_\mathrm{e} = \val{thetaUp}$.
  \item\label{st:state} \emph{The state at the face} \eqref{eq:position},
    \xeqref{first_bend}{eq:momentum}. The ordinate is $y_\mathrm{e} = \val{yUp}$, inside
    $15 \le y \le 19$, and the momentum is $(\val{pxUp}, \val{pyUp})$, at
    $\val{angUp}^\circ$ to the launch. The kinetic energy, from
    \xeqref{foundations}{eq:energy}, is $\val{TUp}$.
  \item\label{st:time} \emph{The time}, Simpson's rule with $\val{nSimp}$ panels
    (§\ref{sec:time}): $t_\mathrm{e} = \val{tUp}$, uncertain by $\val{terrUp}$, and far
    below $t_{\max}$.
  \item\label{st:lower} \emph{The lower shot}, by the same steps. Its constants and its
    result:
    \begin{center}\small\val{orbits}\end{center}
    The second row is the lower shot. The root is $\theta_\mathrm{e} = \val{thetaDown}$,
    the ordinate $\val{yDown}$ lies inside $1 \le y \le 5$, the time is
    $\val{tDown}$ (uncertain by $\val{terrDown}$), and the momentum is
    $(\val{pxDown}, \val{pyDown})$, at $\val{angDown}^\circ$. The closest approach is
    $r = \val{rminDown}$ at $\theta = \val{thMinDown}$, reproduced by the same sampling to
    $\val{rminGapDown}$, and the kinetic energy is $\val{TDown}$.
  \item\label{st:radiation} \emph{The radiated energy} (§\ref{sec:time}):
    $\val{EradUp}$ and $\val{EradDown}$, with relative uncertainties $\val{radErrUp}$ and
    $\val{radErrDown}$.
\end{procedure}

\begin{result}[The reference flights]
Both particles enter their detectors through the face $x = \val{xFace}$. The upper
particle arrives at $t = \val{tUp}$ with ordinate $\val{yUp}$, just inside the edge
$y = 15$, and momentum $(\val{pxUp}, \val{pyUp})$. The lower particle arrives at
$t = \val{tDown}$ with ordinate $\val{yDown}$ and momentum $(\val{pxDown}, \val{pyDown})$.
Both times lie below $t_{\max}$. The level is solved.
\end{result}

Steps \ref{st:launch}--\ref{st:state} and \ref{st:lower} use arithmetic, $\sin$, $\cos$ and
$\arctan$, and the short Newton table. The two quadratures are specified in
§\ref{sec:time}; their errors sit far below $10^{-9}$.

\section{Comparison with the game}\label{sec:comparison}

\comparisontable

The orbit is exact, and the quadratures are uncertain by $\val{terrUp}$ and
$\val{terrDown}$ in the time and by the relative amounts $\val{radErrUp}$ and
$\val{radErrDown}$ in the radiated energy. The tolerance on a state is $10^{-9}$, ten
times the accuracy of a verified flight (\xsecref{foundations}{sec:verification}); the
tolerance on a radiated energy is $10^{-4}$ relative. Both were set from those estimates.

Both reference integrations stop on the face $x = \val{xFace}$, so those abscissae agree
exactly. The largest absolute difference in the table is $\val{cmpHardUpXDiff}$, the
abscissa where the stronger charge's flights leave the world (§\ref{sec:others}); the two
flights share it. On the reference, the largest is $\val{cmpUpYDiff}$, the upper arrival
ordinate. All of them sit at the game's accuracy, above the quadrature errors. The radiated
energies agree to the relative differences $\val{cmpUpRadDiff}$ and $\val{cmpDownRadDiff}$.

\section{Other solutions}\label{sec:others}

Two further placements, each flown by the same orbit and compared above. Both charges lie
on $y = 10$, so §\ref{sec:mirror} applies: equal times, opposite angles, and ordinates that
sum to $20$.

\begin{enumerate}
  \item $Q = \val{SymQ}$ at $\val{SymNode}$, so $\kappa = \val{SymK}$, the reference
    magnitude on the mirror line. Both shots arrive. Each takes $t = \val{SymUpT}$
    (quadrature $\val{SymUpErr}$). The ordinates are $\val{SymUpY}$ and $\val{SymDownY}$,
    and the directions are $\pm\val{SymUpAng}^\circ$. The detectors are mirrors, so one
    arrival forces the other. The reference charge, one cell lower, is the unsymmetric way
    to do the same job: the narrower offset turns more (\eqref{eq:chi}).
  \item $Q = \val{HardQ}$ at $\val{HardNode}$, so $\kappa = \val{HardK}$, twice the
    coupling and two cells closer to the launches. Both shots leave the world, at
    $t = \val{HardUpT}$, through $y = \val{HardUpY}$ and $y = \val{HardDownY}$, at
    $x = \val{HardUpX}$ and at $\pm\val{HardUpAng}^\circ$. The face $x = \val{xFace}$ is
    never reached. The turn has become large enough to hit the ceiling and the floor.
\end{enumerate}

The reference works because the two offsets differ by two cells. \eqref{eq:chi} then turns
the lower ray about twice as far as the upper one, and the detectors are placed one above
the launch line and one below it. A charge on the mirror line also works, with one common
angle. A stronger charge on that line turns both rays into the edges of the world.

\section{What the game computes}\label{sec:engine}

The game integrates \eqref{eq:motion} for each shot separately, with its eighth-order
Runge--Kutta method and step-size control at the tolerances $10^{-10}$ and $10^{-12}$
(\physref{5}). It finds each crossing on its interpolant (\physref{6}) and verifies the
outcome by comparing its two runs (\physref{7}). The model is this document's, including
\ref{as:radiation} (\physref{8}). The scattering formula is not part of what the game
computes. The agreement in §\ref{sec:comparison} is between the orbit
\eqref{eq:orbit} and that integration.

\section{Further reading}\label{sec:reading}

\begin{itemize}
  \item \levelref{first_bend}: the orbit \eqref{eq:orbit}, the time element and the
    momentum. There the launch is an apsis; here $u'_0 \ne 0$, and the time is the
    quadrature of §\ref{sec:time}.
  \item \levelref{geiger_marsden}: the scattering angle \eqref{eq:chi}, and the three ways
    it misses a flight that does not come from infinity.
  \item \xsecref{foundations}{sec:board}: several shots form one puzzle, solved only when
    every shot arrives. This is the first level that asks for that.
\end{itemize}

\levelfileappendix

\end{document}
